%! Author = Javier Mérida
%! Date = 2/12/2023

% Preamble
\documentclass[12pt]{report}
\parskip=\baselineskip
\setcounter{secnumdepth}{4}
% Packages
\usepackage[utf8]{inputenc}
\usepackage{graphicx}
\usepackage[smartEllipses]{markdown}
\usepackage{setspace}
\usepackage[table,xcdraw]{xcolor}
\usepackage{float}
\singlespacing
% Graphics
\graphicspath{ {images/} }
\floatplacement{figure}{H}
% title
\title{
        {Static Site Generator}\\
    {\large La Salle Campus Barcelona}\\
    {\includegraphics{lasalle-logo}}
}
\author{Javier Mérida}
\date{12 02 2023}

% Document
\begin{document}

    \maketitle


    \chapter{Trying out popular SSGs}\label{ch:trying-out-popular-ssgs}

    The following SSGs has been chosen based on their popularity in order
    to be tried to gather information on what the basic features for SSG
    are.
    It's important to notice that these are being tried in a Linux environment
    (Linux inside Windows, thanks to WSL), installing them from scratch, using vim
    as text editor to minimize external tools assistance to have an unbiased appreciation
    of each one of them.

    \begin{itemize}
        \item NextJS (https://nextjs.org/)
        \item Hugo
        \item VuePress
        \item Eleventy
        \item Astro
    \end{itemize}

    % Trying NextJS
    \markdownInput{chapters/testing-ssgs/nextjs.md}

    % Trying Hugo
    \markdownInput{chapters/testing-ssgs/hugo.md}

    % Trying VuePress
    \markdownInput{chapters/testing-ssgs/vuepress.md}

    % Trying Eleventy
    \markdownInput{chapters/testing-ssgs/eleventy.md}

    % Trying Astro
    \markdownInput{chapters/testing-ssgs/astro.md}


    \section{Comparison Table}\label{sec:comparison-table}


    The following table outlines the features that has been considered as most important to take into account when
    considering static site generator.
    Note that the purpose of this document is not to benchmark any of these features,
    thus there is not a workload test to compare the performance of the explored SSGs. Instead, this is a simple
    exploration on what features are stand out in order to consider as requisites to build a new SSG.


    Also, it is important to note that the lack or the presence of any of these characteristic do not make any SSG better
    than others, as each one of them are made for a specific purpose and shine in terms of addressing the obstacle
    they are meant to.
    In fact, this section must be taken as an exploration of requisites to be taken into account.


    % Please add the following required packages to your document preamble:
% \usepackage[table,xcdraw]{xcolor}
% If you use beamer only pass "xcolor=table" option, i.e. \documentclass[xcolor=table]{beamer}
    \begin{table}[H]
        \begin{tabular}{|
                >{\columncolor[HTML]{EFEFEF}}c |c|c|c|c|c|}
            \hline
            \cellcolor[HTML]{C0C0C0}      & \cellcolor[HTML]{EFEFEF}NextJS & \cellcolor[HTML]{EFEFEF}Hugo & \cellcolor[HTML]{EFEFEF}VuePress & \cellcolor[HTML]{EFEFEF}Eleventy & \cellcolor[HTML]{EFEFEF}Astro \\ \hline
            \textbf{Easy to use}          & Yes                            & Yes                          & No                               & Yes                              & Yes                           \\ \hline
            \textbf{Flexible}             & No                             & No                           & No                               & Yes                              & Yes                           \\ \hline
            \textbf{Strict structure}     & No                             & Yes                          & No                               & No                               & No                            \\ \hline
            \textbf{Template system}      & No                             & Yes                          & No                               & Yes                              & Yes                           \\ \hline
            \textbf{Themes}               & No                             & Yes                          & No                               & Yes                              & Yes                           \\ \hline
            \textbf{Performance oriented} & Yes                            & Yes                          & NA                               & Yes                              & Yes                           \\ \hline
            \textbf{Framework dependant}  & Yes                            & No                           & Yes                              & No                               & No                            \\ \hline
            \textbf{Easy to install}      & Yes*                           & Yes                          & Yes*                             & Yes*                             & Yes*                          \\ \hline
            \textbf{File-based routing}   & Yes                            & Yes                          & NA                               & Yes                              & Yes                           \\ \hline
        \end{tabular}\label{tab:table}
    \end{table}


    Note*: The installation difficulty is strickly dependant of NodeJS and how `npm` addresses any missing dependency,
    which can be (sometimes) a painful issue to solve.

    % SRS
    \chapter{Software requirements specifications}\label{ch:software-requirements-specifications}

    % SRS Introduction.
    \markdownInput{chapters/requirements/introduction.md}

    % SRS Overall description
    \markdownInput{chapters/requirements/description.md}

    % SRS Specific requirements.
    \markdownInput{chapters/requirements/requirements.md}

    \setkeys{Gin}{width=1\linewidth}
    % Software design
    \markdownInput{chapters/software-design/design.md}

    % Used technologies
    \markdownInput{chapters/used-tech/technologies.md}

    \setkeys{Gin}{width=0.7\linewidth}
    % Implementation
    \markdownInput{chapters/implementation/implementation.md}

\end{document}
